实例图像目标导航（IIN）任务是具身视觉领域的核心任务之一，它要求智能体在未知环境中，定位目标图像中所包含的特定物体实例。
该任务的核心挑战在于：智能体需要在不同视角下识别物体，定位，同时还需要排除干扰。
在既有方法中，传统的鸟瞰图（BEV）方法存在三维几何信息缺失、实例纹理特征丢失等问题，而基于3D高斯（3DGS）的导航方法语义存在着标签与几何纹理参数脱节的问题。

针对上述问题，本文提出了一种基于语义增强3D 高斯泼溅的模块化视觉导航算法框架\textbf{GSGuide}。
该框架设计并实现了三阶段模块化导航策略，包括\textbf{环境探索与语义地图构建，目标定位与相机位姿对齐，以及智能体规划路径与导航}。
在目标定位与相机位姿估计阶段，提出了图像匹配初始粗定位以及图像扭曲（Warping）渲染精修的高效定位策略。
同时，该框架还通过为每个高斯引入了可训练语义向量，构建了同时编码了几何、纹理与细粒度语义信息的统一地图表示方案。
所提框架使得智能体能够精准地识别和定位目标物体，并能导航至目标所在位置。


% 同时，本文设计了包含交叉熵损失、Dice 损失和 Focal IoU 损失的语义渲染监督机制，实现了语义标签与高斯参数的端到端联合优化；
% 提出 "初始粗定位 + 渲染精修" 的两阶段目标定位策略，利用 DISK+LightGlue 特征匹配获取初始位姿，再通过 Warping Loss 实现目标位姿的亚像素级优化。
基于Habitat 仿真平台和具有挑战性的HM3D 数据集，论文对所提方法完成了仿真实验。
实验结果表明，GSGuide的导航成功率从原始GaussNav 的0.725 提升至0.778，SPL指标达到了0.585，同时在所有语义分割损失指标上均取得了显著改善，验证了所提方法在实例级目标识别与定位方面的有效性，不仅在理论上丰富了具身视觉导航的地图表示方法与目标定位技术，也在实践中为服务机器人的物品查找、室内巡检等实际应用提供了可行的技术方案。